\documentclass[12pt,letterpaper]{article}

\usepackage[margin=1in]{geometry}
\usepackage[T1]{fontenc}
\usepackage{amsmath}
\usepackage{amssymb}
\usepackage{newtxtext,newtxmath}
\usepackage{setspace}
\usepackage{caption}
\usepackage{graphicx}
\usepackage{booktabs,array}
\usepackage{float}
\usepackage{placeins}
\usepackage{titlesec}
\usepackage[hidelinks]{hyperref}

% Match the blueprint's body-size, bold, unnumbered section headings.
\titleformat{\section}{\normalfont\normalsize\bfseries}{}{0pt}{}
\titlespacing*{\section}{0pt}{\baselineskip}{0pt}
\setlength{\parindent}{0pt}
\setlength{\parskip}{0pt}

\begin{document}

\hypersetup{pageanchor=false}
\begin{titlepage}
    \centering
    \singlespacing
    \renewcommand{\thefootnote}{\fnsymbol{footnote}}
    \vspace*{1.5in}

    {\Large\bfseries Does higher education complement or\\
    substitute for family background?\footnote{\normalfont This paper was prepared as a class assignment for Econ 4331W Development Economics.}\par}
    \vspace{\baselineskip}
    Ellie Hou\par
    November 2026\par

    \vfill
\end{titlepage}
\setcounter{footnote}{0}
\hypersetup{pageanchor=true}

\singlespacing
\raggedright

\section*{INTRODUCTION}
My empirical question is: ``Is the adjusted annual earnings premium associated with
a bachelor's degree, relative to a high-school diploma, larger for respondents
from lower-income families?'' This comparison describes how education-related
earnings differences vary by family income and is relevant to debates about higher
education and intergenerational mobility. \newline

Using data from the National Longitudinal Survey of Youth 1997, this paper examines differences
in the college earnings premium by family background, controlling for adolescent academic
preparation and other demographic characteristics. I define financial disadvantage using 
baseline household income, which measures resources available before college. 
I distinguish completed degree levels so that the bachelor's premium is measured separately 
from the graduate/professional premium, using high-school diploma holders as the comparison group.\newline

The descriptive statistics show that a smaller share of the lower-income group completed
a bachelor's degree or higher by 2015. The preliminary empirical results show that the 
adjusted bachelor's premium is smaller for respondents from lower-income families, a pattern 
consistent with complementarity between family resources and higher education.
These regression estimates describe conditional associations and do not establish
causal complementarity, substitution, or effects on intergenerational mobility.\newline

Further research on expected earnings, college costs, financing constraints, 
and enrollment and completion decisions could help explain these patterns and guide a 
research design better suited to identifying causal effects, eventually helping to inform
policy decisions and support a more equitable higher education system in which students
with passion and talent can thrive alongside their wealthier peers.

\section*{LITERATURE REVIEW}

Mincer (1974) links earnings to investments in schooling and work experience,
providing a human-capital framework for examining earnings differences across
educational attainment. My specification draws on this framework but compares
completed degree categories using annual earnings. \newline

Becker and Tomes (1986) model how parental investment decisions and opportunities
transmit economic advantage across generations. De Nardi (2015) reviews related
mechanisms in the wealth-inequality literature, including human-capital transfers
and bequests. These contributions motivate considering family resources alongside
education. \newline

More directly, Chetty et al. (2020) document parental-income segregation across
colleges and examine adult earnings by students' family-income backgrounds.
Chetty, Deming, and Friedman (2026) study admission to highly selective private
colleges and its consequences for later outcomes. Their need-affirmative admissions
exercise uses the bottom parental-income quintile as a policy threshold, following
Chetty et al. (2020). I adapt that threshold to the weighted income distribution in
my NLSY97 sample. \newline

Using a quantitative model of college choice, Hendricks, Koreshkova, and Leukhina
(2026) predict earnings and welfare gains from expanding capacity at selective
public universities, partly by admitting high-ability students previously rationed
out. These model-based policy predictions and the selective-college studies
highlight the potential importance of access and institutional differences
and motivate a follow-up analysis of this paper.\newline

Together, these studies motivate examining whether the bachelor's-versus-high-school 
earnings premium varies by family income, with college selectivity remaining a possible
 explanation for future research.

\section*{HYPOTHESIS}
My primary hypothesis is that the adjusted annual earnings premium associated with
a bachelor's degree only, relative to a high-school diploma, is larger for respondents
from lower-income families. Degree status is measured by 2015, and lower income
means belonging to the weighted bottom quintile of 1996 household income in the sample. 
In the model below, the premium is $\beta_{\mathrm{BA}}$ above the
income cutoff and $\beta_{\mathrm{BA}}+\delta_{\mathrm{BA}}$ at or below it.
My directional prediction is therefore $\delta_{\mathrm{BA}}>0$,
conditional on ASVAB percentile, sex, age, race/ethnicity, and earnings-year effects.\newline

One possible mechanism is that the skills and credentials associated with a
bachelor's degree substitute for advantages otherwise provided by family resources, 
though it can't be established by the estimation in this paper alone. \newline

My secondary hypothesis is that the lower-income group has a smaller share of
respondents who completed a bachelor's degree or higher by 2015. Family resources
and earlier educational investments may affect the ability to finance and complete
college, consistent with the parental-investment framework of Becker and Tomes
(1986). I evaluate this prediction using weighted completion shares within the
selected sample. This descriptive BA-or-higher measure includes graduate and
professional degrees. My primary regression, in contrast, treats bachelor's-only and
graduate/professional degree holders separately.\newline

\section*{DATA}
\textbf{Data source and sample}\\
I use the National Longitudinal Survey of Youth 1997 (NLSY97), which follows a
U.S. cohort born between 1980 and 1984. The analysis uses baseline household
income in 1996, the highest degree earned
by 2015, annual employee earnings and the corresponding income-receipt screens,
the adolescent ASVAB score, sex, birth year, and race/ethnicity. I convert nominal dollars into constant 2025 dollars using the annual average CPI-U,
multiplying each value by $CPI_{2025}/CPI_t$.\footnote{The supplied CPIAUCNS series
has no October 2025 observation, consistent with
\href{https://www.bls.gov/news.release/archives/cpi_12182025.htm}{BLS documentation}.
I use the mean of the 11 available months, 321.943.}\newline

The balanced panel requires observed baseline household income, degree status,
ASVAB, and demographic controls, together with usable employee-earnings measures
for 2016, 2018, and 2020. It contains 3,171 people and 9,513 person-year observations.
All respondents have a positive official custom weight for this selected panel;
the same person weight is used in each earnings year. I retain reported zero earnings and
recover a skipped wage amount as zero only when the income-receipt screening
questions explicitly establish no wage income. Unresolved skips and other missing
amounts are excluded. This recovers 1,532 zero wage amounts in the selected panel.
Genuine negative baseline household income is retained. Family background is
defined by baseline household income.\newline

\textbf{Dependent variable $Y_{it}$}\\
$Y_{it}$ measures the respondent's annual pretax income from wages, salaries,
commissions, and tips across all jobs in 2016, 2018, and 2020, expressed in 2025
US dollars. The source is question YINC-1700 in the 2017, 2019, and 2021 surveys
(variables U2857200, U4282300, and U5753500), each referring to the preceding
calendar year. Its topcoding replaced values in the highest 2\% with 
their group's mean within each survey's valid responses. This limits the
 observed detail in the upper earnings tail.\newline

Annual earnings reflect both compensation rates and time spent in paid work,
including differences in weeks worked and hours worked per week. Consequently,
the estimated degree premiums describe differences in yearly employee earnings;
they cannot isolate an hourly wage premium. I use earnings in levels and include
reported and screening-confirmed zeros, so the estimates are dollar differences
that also reflect whether respondents receive any wage income. Zero employee earnings
do not necessarily imply unemployment, since a respondent may have self-employment
income.\newline

\textbf{Education categories $D_{ik}$}\\
I classify the highest degree earned by the 2015 interview (U0009400) into five mutually exclusive
categories: no high-school diploma or GED only; high-school diploma; associate
degree; bachelor's degree only; and graduate or professional degree. High-school
diploma is the reference category. The source codes map to these groups as
0--1, 2, 3, 4, and 5--7, respectively. This separates associate, bachelor's, and
graduate credentials, while retaining pooled categories for no diploma/GED and
graduate/professional degrees. High-school diploma holders can include people who
attended college without completing a higher degree. I classify respondents by their
highest degree at the 2015 interview and hold that classification fixed across the
three earnings years. A separate descriptive completion indicator
equals one for respondents with a bachelor's or any graduate/professional degree (source codes 4--7).\newline

\textbf{Lower-income background $L_i$}\\
$L_i=1$ if 1996 household income is at or below its weighted 20th-percentile cutoff,
\$21,613 in nominal 1996 dollars. I calculate this cutoff within the selected panel using
\href{https://www.nlsinfo.org/content/weights/nlsy97}{custom NLSY97 weights}, counting
each person once. The cutoff is the smallest observed income at which cumulative
weight reaches 20\%; all ties at the cutoff are included. The lower-income group
contains 837 people and represents 20.03\% of weighted respondents. The comparison
group contains 2,334 people with incomes above the cutoff.\newline

This is an adaptation of the bottom-quintile threshold in Chetty, Deming, and
Friedman (2026, Section VI.A and Table V), while it does not fully replicate 
their approach. They average parental household income over child ages 12--17 
and construct national birth-cohort ranks, while I use one year's household 
income and pool the 1980--84 birth cohorts within the selected sample. The cutoff 
is calculated based on the reported income in this sample.\newline

\textbf{Control variables $X_{it}$}\\
I measure adolescent academic preparation with the NLSY97-created ASVAB
math--verbal percentile, R9829600
(\texttt{ASVAB\_MATH\_VERBAL\_SCORE\_PCT}). NLS staff construct weighted 
percentile ranks within 20 three-month birth cohorts, combining the two math 
components with twice the percentile of the verbal composite and then ranking 
the aggregate within each cohort (U.S. Bureau of Labor Statistics, n.d.-b). 
The supplied score therefore already accounts for age-group differences through
its birth-cohort normalization.\newline

I divide the stored score by 1,000 to restore its three implied decimal places and
express it in percentile points. I retain the supplied percentile and do not rerank
it within my selected panel. The weights NLS used to construct the score were based
on a different sample and are
distinct from the custom weights used in my earnings regression.\newline

I include ASVAB percentile linearly, alongside a female indicator (male reference),
age at the end of the earnings year, and race/ethnicity indicators for Black,
Hispanic, and mixed non-Hispanic respondents (non-Black/non-Hispanic reference).
Age is the earnings calendar year minus birth year. Earnings-year indicators use
2016 as the reference. These controls accompany the four education indicators,
the lower-income indicator, and their four interactions. The ASVAB score captures skills measured during
adolescence and can reflect prior schooling and family resources. Thus the
comparison conditions on measured preparation; this control does not eliminate
unobserved differences in motivation, skills, or selection into degree completion.\newline

Table~\ref{tab:summary-statistics-current} summarizes the analytic panel by baseline
income group, using the same custom person weights as the regression.

\IfFileExists{tables/summary_statistics_current.tex}{% Snapshot calculated from data/processed/analysis_panel.csv; refresh after changing the analysis.
% Source SHA256: e1832590e1f7343920a575afc5e03e8b71cc229244369c913f30ad3b42e377a4
% SD is sqrt(sum(w*(x-weighted_mean)^2)/sum(w)); person-years only for earnings, zero earnings, and age.
\begin{table}[!htbp]
\centering
\singlespacing
\small
\captionsetup{font=small,labelfont=bf,justification=raggedright,singlelinecheck=false}
\caption{Summary statistics by baseline household income}
\label{tab:summary-statistics-current}
\setlength{\tabcolsep}{4pt}
\renewcommand{\arraystretch}{1.12}
\begin{tabular}{@{}p{0.40\linewidth}*{3}{>{\centering\arraybackslash}p{0.18\linewidth}}@{}}
\toprule
Variable & Full sample & Lower income & Above cutoff \\
\midrule
Annual employee earnings (\$1,000) & 68.6 (72.4) & 43.7 (51.8) & 74.8 (75.4) \\
Household income, 1996 (\$1,000) & 110.1 (88.6) & 24.9 (14.2) & 131.5 (86.6) \\
ASVAB percentile points & 53.8 (28.8) & 35.4 (27.0) & 58.4 (27.3) \\
Age in earnings year & 36.0 (2.2) & 35.9 (2.1) & 36.0 (2.2) \\
Zero annual employee earnings (\%) & 15.4 & 26.7 & 12.5 \\
\addlinespace[4pt]
Female (\%) & 48.7 & 52.6 & 47.7 \\
Black (\%) & 15.3 & 28.8 & 12.0 \\
Hispanic (\%) & 12.8 & 19.9 & 11.0 \\
Mixed non-Hispanic (\%) & 1.1 & 0.8 & 1.2 \\
Non-Black/non-Hispanic (\%) & 70.8 & 50.4 & 75.9 \\
\addlinespace[4pt]
No diploma/GED (\%) & 14.4 & 32.4 & 9.9 \\
High-school diploma (\%) & 39.1 & 45.1 & 37.7 \\
Associate degree (\%) & 9.1 & 8.0 & 9.4 \\
Bachelor's degree only (\%) & 25.3 & 10.1 & 29.1 \\
Graduate/professional degree (\%) & 12.1 & 4.5 & 14.0 \\
Bachelor's degree or higher (\%) & 37.3 & 14.6 & 43.0 \\
\midrule
Respondents, unweighted & 3,171 & 837 & 2,334 \\
\bottomrule
\end{tabular}
\par\smallskip
\begin{minipage}{\linewidth}
\footnotesize\singlespacing
\textit{Notes:} Continuous variables show weighted means with weighted standard
deviations in parentheses; categorical variables show weighted percentages.
All monetary values are expressed in
thousands of 2025 US dollars. Earnings, zero-earnings shares, and age pool
person-years from 2016, 2018, and 2020; other statistics count each respondent once.
Each respondent contributes three observations, covering these three earnings years.
Respondent counts are unweighted. Lower income means household income at or below \$21,613
in nominal 1996 dollars. The BA-or-higher row combines the bachelor's and graduate
categories. Percentages may not sum to 100 because of rounding.
\end{minipage}
\end{table}
}{% Snapshot calculated from data/processed/analysis_panel.csv; refresh after changing the analysis.
% Source SHA256: e1832590e1f7343920a575afc5e03e8b71cc229244369c913f30ad3b42e377a4
% SD is sqrt(sum(w*(x-weighted_mean)^2)/sum(w)); person-years only for earnings, zero earnings, and age.
\begin{table}[!htbp]
\centering
\singlespacing
\small
\captionsetup{font=small,labelfont=bf,justification=raggedright,singlelinecheck=false}
\caption{Summary statistics by baseline household income}
\label{tab:summary-statistics-current}
\setlength{\tabcolsep}{4pt}
\renewcommand{\arraystretch}{1.12}
\begin{tabular}{@{}p{0.40\linewidth}*{3}{>{\centering\arraybackslash}p{0.18\linewidth}}@{}}
\toprule
Variable & Full sample & Lower income & Above cutoff \\
\midrule
Annual employee earnings (\$1,000) & 68.6 (72.4) & 43.7 (51.8) & 74.8 (75.4) \\
Household income, 1996 (\$1,000) & 110.1 (88.6) & 24.9 (14.2) & 131.5 (86.6) \\
ASVAB percentile points & 53.8 (28.8) & 35.4 (27.0) & 58.4 (27.3) \\
Age in earnings year & 36.0 (2.2) & 35.9 (2.1) & 36.0 (2.2) \\
Zero annual employee earnings (\%) & 15.4 & 26.7 & 12.5 \\
\addlinespace[4pt]
Female (\%) & 48.7 & 52.6 & 47.7 \\
Black (\%) & 15.3 & 28.8 & 12.0 \\
Hispanic (\%) & 12.8 & 19.9 & 11.0 \\
Mixed non-Hispanic (\%) & 1.1 & 0.8 & 1.2 \\
Non-Black/non-Hispanic (\%) & 70.8 & 50.4 & 75.9 \\
\addlinespace[4pt]
No diploma/GED (\%) & 14.4 & 32.4 & 9.9 \\
High-school diploma (\%) & 39.1 & 45.1 & 37.7 \\
Associate degree (\%) & 9.1 & 8.0 & 9.4 \\
Bachelor's degree only (\%) & 25.3 & 10.1 & 29.1 \\
Graduate/professional degree (\%) & 12.1 & 4.5 & 14.0 \\
Bachelor's degree or higher (\%) & 37.3 & 14.6 & 43.0 \\
\midrule
Respondents, unweighted & 3,171 & 837 & 2,334 \\
\bottomrule
\end{tabular}
\par\smallskip
\begin{minipage}{\linewidth}
\footnotesize\singlespacing
\textit{Notes:} Continuous variables show weighted means with weighted standard
deviations in parentheses; categorical variables show weighted percentages.
All monetary values are expressed in
thousands of 2025 US dollars. Earnings, zero-earnings shares, and age pool
person-years from 2016, 2018, and 2020; other statistics count each respondent once.
Each respondent contributes three observations, covering these three earnings years.
Respondent counts are unweighted. Lower income means household income at or below \$21,613
in nominal 1996 dollars. The BA-or-higher row combines the bachelor's and graduate
categories. Percentages may not sum to 100 because of rounding.
\end{minipage}
\end{table}
}
\FloatBarrier

\section*{METHODOLOGY}

The primary model is
\[
\begin{aligned}
    Y_{it}={}&\alpha+\gamma L_i+\sum_{k\in\mathcal K}\beta_k D_{ik}
              +\sum_{k\in\mathcal K}\delta_k(D_{ik}\times L_i)
             &+X'_{it}\theta+\lambda_t+\epsilon_{it}
\end{aligned}
\]
where $\mathcal K$ contains the four education categories other than high-school
diploma, and $D_{ik}$ indicates membership in category $k$. $\beta_k$ is the adjusted
earnings difference between category $k$ and high-school diploma holders among
respondents above the income cutoff; $\beta_k+\delta_k$ is the corresponding
difference for the lower-income group. Thus $\delta_{\mathrm{BA}}$ compares the
bachelor's-only premium across income groups. $\lambda_t$ controls for earnings-year
effects, and $X_{it}$ includes ASVAB percentile, sex, age, and race/ethnicity.\newline

I estimate weighted pooled OLS with year effects, applying the same custom person
weight in each year, and cluster standard errors by person using the CR1 covariance
estimator. Two-sided tests and 95\% confidence intervals use the Student-$t$
distribution with 3,170 degrees of freedom. Confidence intervals for
$\beta_k+\delta_k$ use the full coefficient covariance. The income-only 20\%
definition was chosen from the literature before fitting this revised model. I also 
report a joint test that all four interactions are zero. \newline

\section*{RESULTS}
Using the custom weights, 14.6\% of the lower-income group completed a bachelor's
degree or higher, compared with 43.0\% of the other group. This descriptive comparison
is consistent with my secondary hypothesis. These results describe weighted
associations within the selected balanced panel.\newline

Table~\ref{tab:education-premia} reports adjusted earnings differences by education
and family income from the same weighted pooled regression. Each education group
is compared with high-school diploma holders in the same income group. The last
column compares these within-group premiums across income groups.\newline

\IfFileExists{tables/education_premia.tex}{% Generated from outputs/regression/education_premia.csv; update after rerunning the model.
% Source SHA256: 43894ebbf0dc1db33b118281b8a668ea4dbd15839d9ae77de457fb1c2bffb7e9
\begin{table}[!htbp]
\centering
\singlespacing
\small
\captionsetup{font=small,labelfont=bf,justification=raggedright,singlelinecheck=false}
\caption{Adjusted annual earnings differences relative to high-school diploma holders}
\label{tab:education-premia}
\setlength{\tabcolsep}{4pt}
\renewcommand{\arraystretch}{1.25}
\begin{tabular}{@{}p{0.34\linewidth}*{3}{>{\centering\arraybackslash}p{0.20\linewidth}}@{}}
\toprule
Highest degree by 2015 & Above income cutoff & Lower income & Difference: lower minus above \\
 & $\widehat\beta_k$ & $\widehat\beta_k+\widehat\delta_k$ & $\widehat\delta_k$ \\
\midrule
No diploma/GED & \shortstack{$-9{,}418^{***}$\\$(2{,}819)$} & \shortstack{$-10{,}939^{***}$\\$(3{,}976)$} & \shortstack{$-1{,}521$\\$(4{,}818)$} \\[3pt]
Associate degree & \shortstack{$15{,}473^{***}$\\$(5{,}066)$} & \shortstack{$7{,}048$\\$(5{,}384)$} & \shortstack{$-8{,}426$\\$(7{,}369)$} \\[3pt]
Bachelor's degree only & \shortstack{$32{,}309^{***}$\\$(3{,}507)$} & \shortstack{$14{,}001^{**}$\\$(5{,}788)$} & \shortstack{$-18{,}308^{***}$\\$(6{,}618)$} \\[3pt]
Graduate/professional & \shortstack{$59{,}237^{***}$\\$(5{,}512)$} & \shortstack{$75{,}540^{***}$\\$(18{,}642)$} & \shortstack{$16{,}303$\\$(19{,}355)$} \\[3pt]
\midrule
Person-year observations & \multicolumn{3}{c}{9,513} \\
Respondents & \multicolumn{3}{c}{3,171} \\
Weighted $R^2$ & \multicolumn{3}{c}{0.221} \\
\bottomrule
\end{tabular}
\par\smallskip
\begin{minipage}{\linewidth}
\footnotesize\singlespacing
\textit{Notes:} Weighted pooled OLS estimates in 2025 US dollars. Each education
category is compared with high-school diploma holders within the same income
group; the final column reports lower-income minus above-cutoff premiums.
Lower income denotes the weighted bottom quintile of baseline household income
in the analytic sample. Respondent-clustered standard errors appear in parentheses.
Controls are described in Methodology. Stars indicate two-sided, unadjusted
significance levels: $^{*}p<0.10$, $^{**}p<0.05$, $^{***}p<0.01$.
\end{minipage}
\end{table}
}{% Generated from outputs/regression/education_premia.csv; update after rerunning the model.
% Source SHA256: 43894ebbf0dc1db33b118281b8a668ea4dbd15839d9ae77de457fb1c2bffb7e9
\begin{table}[!htbp]
\centering
\singlespacing
\small
\captionsetup{font=small,labelfont=bf,justification=raggedright,singlelinecheck=false}
\caption{Adjusted annual earnings differences relative to high-school diploma holders}
\label{tab:education-premia}
\setlength{\tabcolsep}{4pt}
\renewcommand{\arraystretch}{1.25}
\begin{tabular}{@{}p{0.34\linewidth}*{3}{>{\centering\arraybackslash}p{0.20\linewidth}}@{}}
\toprule
Highest degree by 2015 & Above income cutoff & Lower income & Difference: lower minus above \\
 & $\widehat\beta_k$ & $\widehat\beta_k+\widehat\delta_k$ & $\widehat\delta_k$ \\
\midrule
No diploma/GED & \shortstack{$-9{,}418^{***}$\\$(2{,}819)$} & \shortstack{$-10{,}939^{***}$\\$(3{,}976)$} & \shortstack{$-1{,}521$\\$(4{,}818)$} \\[3pt]
Associate degree & \shortstack{$15{,}473^{***}$\\$(5{,}066)$} & \shortstack{$7{,}048$\\$(5{,}384)$} & \shortstack{$-8{,}426$\\$(7{,}369)$} \\[3pt]
Bachelor's degree only & \shortstack{$32{,}309^{***}$\\$(3{,}507)$} & \shortstack{$14{,}001^{**}$\\$(5{,}788)$} & \shortstack{$-18{,}308^{***}$\\$(6{,}618)$} \\[3pt]
Graduate/professional & \shortstack{$59{,}237^{***}$\\$(5{,}512)$} & \shortstack{$75{,}540^{***}$\\$(18{,}642)$} & \shortstack{$16{,}303$\\$(19{,}355)$} \\[3pt]
\midrule
Person-year observations & \multicolumn{3}{c}{9,513} \\
Respondents & \multicolumn{3}{c}{3,171} \\
Weighted $R^2$ & \multicolumn{3}{c}{0.221} \\
\bottomrule
\end{tabular}
\par\smallskip
\begin{minipage}{\linewidth}
\footnotesize\singlespacing
\textit{Notes:} Weighted pooled OLS estimates in 2025 US dollars. Each education
category is compared with high-school diploma holders within the same income
group; the final column reports lower-income minus above-cutoff premiums.
Lower income denotes the weighted bottom quintile of baseline household income
in the analytic sample. Respondent-clustered standard errors appear in parentheses.
Controls are described in Methodology. Stars indicate two-sided, unadjusted
significance levels: $^{*}p<0.10$, $^{**}p<0.05$, $^{***}p<0.01$.
\end{minipage}
\end{table}
}

The adjusted bachelor's-only versus high-school earnings premium is \$14,001 for
the lower-income group and \$32,309 for respondents above the cutoff, in 2025 dollars.
The estimated difference is $\hat\delta_{\mathrm{BA}}=-\$18,308$
(clustered standard error \$6,618; $p=0.0057$; 95\% confidence interval:
$-\$31,285$ to $-\$5,332$). This estimate runs against my hypothesis of a larger
bachelor's premium for lower-income students. Both groups have positive adjusted
bachelor's premiums; the negative interaction indicates that the lower-income
group's premium is smaller. In addition, because the outcome is annual earnings, 
this difference can reflect differences in pay rates, hours, and weeks worked as well.\newline

The other three education interactions are not statistically distinguishable from
zero at the 5\% level in the pointwise tests. In particular, the graduate/professional
interaction is \$16,303 with a clustered standard error of \$19,355. Only 30
lower-income respondents hold graduate or professional degrees, compared with
312 above the income cutoff; the data do not support a precise comparison for
that category. The joint test of all four education interactions yields
$F(4,3170)=2.533$ ($p=0.038$), which rejects their joint equality to zero
at the 5\% level.\newline

Figure~\ref{fig:education-premia} illustrates these comparisons. Panel A shows
positive bachelor's premiums for both income groups, while the bachelor's
premium difference has a confidence interval entirely below zero in Panel B.
The wide graduate/professional interval shows the limited precision of that
comparison. Table~\ref{tab:robustness-checks} in the appendix reports two
robustness checks.

\begin{figure}[H]
\centering
\singlespacing
\captionsetup{font=small,labelfont=bf,justification=raggedright,singlelinecheck=false}
\IfFileExists{figures/education_premia.pdf}{\includegraphics[width=\linewidth]{figures/education_premia.pdf}}{\includegraphics[width=\linewidth]{reports/figures/education_premia.pdf}}
\caption{Adjusted education earnings premiums by family income}
\label{fig:education-premia}
\begin{minipage}{\linewidth}
\footnotesize\singlespacing
\textit{Notes:} Estimates are from the regression in Table~\ref{tab:education-premia},
expressed in thousands of 2025 US dollars. Panel A shows premiums relative to
high-school diploma holders within each income group; Panel B shows lower-income
minus above-cutoff premiums. Lower income denotes the weighted bottom quintile
of baseline household income in the analytic sample. Bars indicate pointwise 95\%
confidence intervals with respondent-clustered standard errors. Panels use
different horizontal scales.
\end{minipage}
\end{figure}
\FloatBarrier

\section*{EXTENSIONS}
Degree attainment in our sample is not randomly assigned, so differences in the 
estimated earnings premium may reflect selection into college enrollment and completion. 
Lower-income students may attend less selective institutions, while financial constraints 
and expected returns may influence which students ultimately obtain a degree. \newline

One possibility is that lower-income degree holders are more strongly selected on actual
earnings gains than graduates from other backgrounds. If other selection differences do
not offset this pattern, the observed premium gap in favor of higher-income respondents
could understate differences in average causal returns. However, higher expected net
returns do not necessarily imply larger realized annual earnings gains: college costs,
forgone earnings, risk, and expectations also affect decisions. The current regression
cannot establish the direction of selection bias.\newline

Information on expected earnings, tuition and financial aid, financing constraints, and
enrollment and completion histories would help investigate these decisions and improve
comparability. College selectivity should be studied as a possible pathway through which
educational opportunities affect earnings. Establishing causal effects would additionally 
require a research design that addresses unobserved selection into degree completion. 
Information on hours and weeks worked could also help distinguish differences in pay 
from differences in time worked.

\section*{REFERENCES}
\begin{list}{}{\setlength{\leftmargin}{0.3in}
              \setlength{\itemindent}{-0.3in}
              \setlength{\itemsep}{0.6\baselineskip}
              \setlength{\parsep}{0pt}}
\item Becker, Gary S., and Nigel Tomes. 1986. ``Human Capital and the Rise and Fall of Families.'' \textit{Journal of Labor Economics} 4 (3, Part 2): S1--S39.
\item Chetty, Raj, David J. Deming, and John N. Friedman. 2026. ``Diversifying Society's Leaders? The Determinants and Causal Effects of Admission to Highly Selective Private Colleges.'' \textit{Quarterly Journal of Economics} 141 (1): 51--145. \href{https://doi.org/10.1093/qje/qjaf050}{doi:10.1093/qje/qjaf050}.
\item Chetty, Raj, John N. Friedman, Emmanuel Saez, Nicholas Turner, and Danny Yagan. 2020. ``Income Segregation and Intergenerational Mobility Across Colleges in the United States.'' \textit{Quarterly Journal of Economics} 135 (3): 1567--1633. \href{https://doi.org/10.1093/qje/qjaa005}{doi:10.1093/qje/qjaa005}.
\item De Nardi, Mariacristina. 2015. ``Quantitative Models of Wealth Inequality: A Survey.'' NBER Working Paper 21106.
\item Hendricks, L., T. Koreshkova, and O. Leukhina. 2026. ``Expanding Access to Selective Colleges.'' Federal Reserve Bank of St. Louis Working Paper 2026-005C. June. \href{https://fedinprint.org/item/fedlwp/102965/103406}{Working-paper record}.
\item Mincer, Jacob A. 1974. \textit{Schooling, Experience, and Earnings}. National Bureau of Economic Research.
\item U.S. Bureau of Labor Statistics. n.d.-a. ``Administration of the CAT-ASVAB.'' \textit{National Longitudinal Surveys, NLSY97 Topical Guide}. \href{https://nlsinfo.org/content/cohorts/nlsy97/topical-guide/education/administration-cat-asvab-0}{Official documentation}. Accessed November 13, 2026.
\item U.S. Bureau of Labor Statistics. n.d.-b. ``Appendix 10: CAT-ASVAB Scores.'' \textit{National Longitudinal Surveys, NLSY97 Codebook Supplement}. \href{https://nlsinfo.org/content/cohorts/nlsy97/other-documentation/codebook-supplement/appendix-10-cat-asvab-scores}{Official documentation}. Accessed November 13, 2026.
\item U.S. Bureau of Labor Statistics. n.d.-c. ``Income.'' \textit{National Longitudinal Surveys, NLSY97 Topical Guide}. \href{https://nlsinfo.org/content/cohorts/nlsy97/topical-guide/income/income}{Official documentation}. See also YINC-1700 entries for the 2017, 2019, and 2021 surveys in the downloaded NLSY97 codebook. Accessed November 13, 2026.
\end{list}

\clearpage
\appendix
\setcounter{table}{0}
\renewcommand{\thetable}{A\arabic{table}}
\renewcommand{\theHtable}{A\arabic{table}}
\section*{APPENDIX: ROBUSTNESS CHECKS}

Table~\ref{tab:robustness-checks} compares the main specification with two
exploratory checks. Each changes one feature while retaining the same 3,171
respondents, official person weights, education categories, and control variables.\newline

\textbf{Excluding 2020.} I use only 2016 and 2018 earnings to assess whether the
result depends on including 2020. The original three-year sample eligibility is
retained, so this check addresses the earnings periods used in estimation rather
than selection from requiring a usable 2020 response.\newline

\textbf{Alternative income cutoff.} I replace the weighted bottom quintile with
the bottom quartile. The resulting cutoff is \$26,000 in nominal 1996 dollars.
This changes the income-group definition while keeping the earnings years and
selected respondents fixed.

\IfFileExists{tables/robustness_checks.tex}{% Snapshot of the main estimate and two exploratory robustness checks.
% outputs/regression/education_premia.csv SHA256: 43894ebbf0dc1db33b118281b8a668ea4dbd15839d9ae77de457fb1c2bffb7e9
% outputs/sensitivity/exclude_2020/regression/education_premia.csv SHA256: e11fad8d10e829941d6697a1b09de52b9d15caffef04702236af6b363c8a55a2
% outputs/sensitivity/income_p25/regression/education_premia.csv SHA256: d4dc211e7002ff57001e34acb60dd778cafde0d0562cbd715b4071e82c8a1dad
\begin{table}[H]
\centering
\singlespacing
\small
\captionsetup{font=small,labelfont=bf,justification=raggedright,singlelinecheck=false}
\caption{Robustness of the bachelor's-only earnings premium}
\label{tab:robustness-checks}
\setlength{\tabcolsep}{3pt}
\renewcommand{\arraystretch}{1.08}
\begin{tabular}{@{}p{0.40\linewidth}*{3}{>{\centering\arraybackslash}p{0.18\linewidth}}@{}}
\toprule
 & Main & Exclude 2020 & 25th percentile \\
\midrule
BA premium: above cutoff & \shortstack{$32{,}309^{***}$\\$(3{,}507)$} & \shortstack{$32{,}523^{***}$\\$(3{,}395)$} & \shortstack{$32{,}437^{***}$\\$(3{,}609)$} \\[2pt]
BA premium: lower income & \shortstack{$14{,}001^{**}$\\$(5{,}788)$} & \shortstack{$11{,}165^{**}$\\$(5{,}494)$} & \shortstack{$19{,}443^{***}$\\$(5{,}719)$} \\[2pt]
Difference: lower minus above & \shortstack{$-18{,}308^{***}$\\$(6{,}618)$} & \shortstack{$-21{,}358^{***}$\\$(6{,}315)$} & \shortstack{$-12{,}994^{**}$\\$(6{,}560)$} \\[2pt]
95\% CI for the difference & \shortstack{$[-31{,}285,$\\$-5{,}332]$} & \shortstack{$[-33{,}739,$\\$-8{,}976]$} & \shortstack{$[-25{,}856,$\\$-131]$} \\
$p$-value for the difference & 0.0057 & 0.0007 & 0.0477 \\
Joint interaction test: $p$-value & 0.0385 & 0.0071 & 0.1857 \\
\midrule
Earnings years & 2016, 2018, 2020 & 2016, 2018 & 2016, 2018, 2020 \\
Lower-income percentile & 20th & 20th & 25th \\
Income cutoff (1996 dollars) & \$21,613 & \$21,613 & \$26,000 \\
Respondents & 3,171 & 3,171 & 3,171 \\
Lower-income respondents & 837 & 837 & 1,004 \\
Person-year observations & 9,513 & 6,342 & 9,513 \\
\bottomrule
\end{tabular}
\par\smallskip
\begin{minipage}{\linewidth}
\footnotesize\singlespacing
\textit{Notes:} Weighted pooled OLS premiums relative to high-school diploma
holders within each income group, in 2025 US dollars. All columns retain the same
respondents, custom weights, education categories, and controls; year effects
correspond to the included years. Respondent-clustered standard errors appear in
parentheses; intervals are pointwise 95\% confidence intervals. The joint test
restricts all four education-by-background interactions to zero. Stars indicate
two-sided, unadjusted significance: $^{*}p<0.10$, $^{**}p<0.05$, $^{***}p<0.01$.
\end{minipage}
\end{table}
}{% Snapshot of the main estimate and two exploratory robustness checks.
% outputs/regression/education_premia.csv SHA256: 43894ebbf0dc1db33b118281b8a668ea4dbd15839d9ae77de457fb1c2bffb7e9
% outputs/sensitivity/exclude_2020/regression/education_premia.csv SHA256: e11fad8d10e829941d6697a1b09de52b9d15caffef04702236af6b363c8a55a2
% outputs/sensitivity/income_p25/regression/education_premia.csv SHA256: d4dc211e7002ff57001e34acb60dd778cafde0d0562cbd715b4071e82c8a1dad
\begin{table}[H]
\centering
\singlespacing
\small
\captionsetup{font=small,labelfont=bf,justification=raggedright,singlelinecheck=false}
\caption{Robustness of the bachelor's-only earnings premium}
\label{tab:robustness-checks}
\setlength{\tabcolsep}{3pt}
\renewcommand{\arraystretch}{1.08}
\begin{tabular}{@{}p{0.40\linewidth}*{3}{>{\centering\arraybackslash}p{0.18\linewidth}}@{}}
\toprule
 & Main & Exclude 2020 & 25th percentile \\
\midrule
BA premium: above cutoff & \shortstack{$32{,}309^{***}$\\$(3{,}507)$} & \shortstack{$32{,}523^{***}$\\$(3{,}395)$} & \shortstack{$32{,}437^{***}$\\$(3{,}609)$} \\[2pt]
BA premium: lower income & \shortstack{$14{,}001^{**}$\\$(5{,}788)$} & \shortstack{$11{,}165^{**}$\\$(5{,}494)$} & \shortstack{$19{,}443^{***}$\\$(5{,}719)$} \\[2pt]
Difference: lower minus above & \shortstack{$-18{,}308^{***}$\\$(6{,}618)$} & \shortstack{$-21{,}358^{***}$\\$(6{,}315)$} & \shortstack{$-12{,}994^{**}$\\$(6{,}560)$} \\[2pt]
95\% CI for the difference & \shortstack{$[-31{,}285,$\\$-5{,}332]$} & \shortstack{$[-33{,}739,$\\$-8{,}976]$} & \shortstack{$[-25{,}856,$\\$-131]$} \\
$p$-value for the difference & 0.0057 & 0.0007 & 0.0477 \\
Joint interaction test: $p$-value & 0.0385 & 0.0071 & 0.1857 \\
\midrule
Earnings years & 2016, 2018, 2020 & 2016, 2018 & 2016, 2018, 2020 \\
Lower-income percentile & 20th & 20th & 25th \\
Income cutoff (1996 dollars) & \$21,613 & \$21,613 & \$26,000 \\
Respondents & 3,171 & 3,171 & 3,171 \\
Lower-income respondents & 837 & 837 & 1,004 \\
Person-year observations & 9,513 & 6,342 & 9,513 \\
\bottomrule
\end{tabular}
\par\smallskip
\begin{minipage}{\linewidth}
\footnotesize\singlespacing
\textit{Notes:} Weighted pooled OLS premiums relative to high-school diploma
holders within each income group, in 2025 US dollars. All columns retain the same
respondents, custom weights, education categories, and controls; year effects
correspond to the included years. Respondent-clustered standard errors appear in
parentheses; intervals are pointwise 95\% confidence intervals. The joint test
restricts all four education-by-background interactions to zero. Stars indicate
two-sided, unadjusted significance: $^{*}p<0.10$, $^{**}p<0.05$, $^{***}p<0.01$.
\end{minipage}
\end{table}
}

The BA premium difference remains negative in both checks. With the
25th-percentile cutoff, the gap falls to \$12,994 and its 95\% confidence interval
nearly reaches zero. The joint test of the four interactions also no longer
rejects at the 5\% level. Thus the direction of the BA association persists,
while its magnitude and the strength of the evidence depend on the income
definition. These checks do not resolve unobserved selection into degree completion.
\end{document}
